\section{Challenges with Nonlinear Dynamics}
\label{sec:discussion-challenges}

Extending formal Lyapunov verification from linear to nonlinear systems introduces theoretical and numerical challenges. 
Unlike the double integrator, the cartpole dynamics involve complex state couplings and trigonometric functions, 
which makes verification harder due to larger relaxation errors. 

Since the DARE solution provides a local Lyapunov approximation around the equilibrium for nonlinear systems and not globally valid, 
its suitability as an initial inductive prior strongly depends on the operating region of the imitation policy. 
When the Lyapunov candidate remains dominated by this quadratic DARE term, the optimization may favor changes to 
the policy that provide a low-cost way to reduce the penalized condition violations. 
Consequently, the optimization may promote large control inputs to reduce 
the Lyapunov decrease condition on mined counterexamples, which are only evaluated pointwise. 
Nevertheless, although the closed-loop state trajectories may exceed the admissible domain in rollouts, 
it is open to question whether the policy remains reasonably close to the expert within the actual training domain $\mathcal{X}_V$.

During formal verification, the linear system benefits from exact bound propagation through its affine dynamics, 
while the nonlinear system suffers from larger relaxation errors. 
Especially the next Lyapunov value $V(f(x,\pi_\theta(x)))$ in \labelcref{eq:lyapunov-decrease} is sensitive to relaxations 
in the policy, because these propagate through the dynamics and the neural Lyapunov candidate. 

Beyond verification, the co-training reveals a limitation of single-step adversarial training sheme. 
As illustrated in \cref{fig:app-cp-cotrain-trajectories}, the co-trained controller exhibits actuator oscillations 
ompared to the smooth imitation baseline. 
Since the adversarial falsifier is restricted to single-step transitions $x_k \mapsto x_{k+1}$, 
the current falsification scheme does not capture multi-step dynamic accumulation.
By generating localized extreme control inputs, the policy can reduce the decrease violations 
at the sampled points, while potentially introducing less smooth control behavior over longer rollouts.